\documentclass[11pt]{article}
\usepackage[T1]{fontenc}
\usepackage{graphicx,hyperref}
\usepackage{titling}

\usepackage{xcolor}
\usepackage{amsmath,amssymb}
\usepackage[normalem]{ulem}
\usepackage{fancyvrb}
\usepackage[most]{tcolorbox}
\usepackage{microtype}

\newtcolorbox{srccode}{
  enhanced, breakable,
  colback=white, colframe=black!20,
  left=4pt, right=4pt, top=2pt, bottom=2pt,
  boxrule=0.4pt,
  arc=0pt, outer arc=0pt
}

\newtcolorbox{outcode}{
  enhanced, breakable,
  colback=white, colframe=black!20,
  left=4pt, right=4pt, top=2pt, bottom=2pt,
  boxrule=0.4pt,
  arc=0pt, outer arc=0pt,
  fontupper=\small
}

\title{Basics}
\posttitle{\par\vskip 0.5em\large \end{center}}


\definecolor{hl1}{RGB}{51,51,51}
\definecolor{hl2}{RGB}{167,29,93}
\definecolor{hl3}{RGB}{0,134,179}
\definecolor{hl4}{RGB}{150,152,150}
\definecolor{hl5}{RGB}{24,54,145}
\begin{document}
\maketitle
\ifdefined\affiliationblock\affiliationblock\fi





This page reviews some of the basics of Quarto and Markdown syntax supported by \emph{Calepin}. It covers document structure, code chunks, figures, math, cross-references, citations, and other features that work out of the box.

\subsection*{Getting started}
\label{getting-started}

\emph{Calepin} renders \texttt{.qmd} documents to HTML, LaTeX, Typst, and Markdown. It executes R code chunks, captures output and plots, resolves citations and cross-references, and wraps everything in a template.

To render this tutorial:




\begin{srccode}
\begin{Verbatim}[commandchars=\\\{\}]
\textcolor{hl1}{calepin tutorial.qmd}
\end{Verbatim}
\end{srccode}

\subsection*{Markdown basics}
\label{markdown-basics}

\emph{Calepin} supports standard CommonMark and GitHub Flavored Markdown. The table below summarizes the basic Markdown notation from the original specification.

\begin{tabular}{ll}
\hline
Element & Syntax \\
\hline
Heading & \texttt{\# H1}, \texttt{\#\# H2}, \ldots{}, \texttt{\#\#\#\#\#\# H6} \\
Paragraph & Blank line between blocks of text \\
Line break & Two trailing spaces or \texttt{\textbackslash{}} at end of line \\
Italic & \texttt{*text*} or \texttt{\_text\_} \\
Bold & \texttt{**text**} or \texttt{\_\_text\_\_} \\
Bold + italic & \texttt{***text***} \\
Blockquote & \texttt{\textgreater{} quoted text} \\
Ordered list & \texttt{1. item} \\
Unordered list & \texttt{- item}, \texttt{* item}, or \texttt{+ item} \\
Inline code & \texttt{`code`} \\
Code block & Indent by 4 spaces or fence with \texttt{```} \\
Horizontal rule & \texttt{-{}-{}-}, \texttt{***}, or \texttt{\_\_\_} \\
Link & \texttt{{[}text{]}(url)} or \texttt{{[}text{]}(url \textquotedbl{}title\textquotedbl{})} \\
Reference link & \texttt{{[}text{]}{[}id{]}} then \texttt{{[}id{]}: url} \\
Image & \texttt{!{[}alt text{]}(path)} \\
Escape & \texttt{\textbackslash{}*}, \texttt{\textbackslash{}{[}}, \texttt{\textbackslash{}\textbackslash{}}, etc. \\
\hline
\end{tabular}

Here is a paragraph about Montreal that illustrates all of these features.


\begin{center}\rule{0.5\linewidth}{\linethickness}\end{center}

\subsubsection*{Montréal}
\label{montréal}

\begin{quote}
\emph{Montréal is a city of contrasts: old-world charm meets modern ambition.}

\end{quote}

Montréal is the largest city in the province of Québec and the \emph{second most populous} in Canada. Founded in 1642, the city sits on an island in the \emph{Saint Lawrence River}. Visitors should explore these landmarks:

\begin{enumerate}
\item Old Montréal: cobblestone streets and the \href{https://www.basiliquenotredame.ca}{Notre-Dame Basilica}


\item Mount Royal: the park designed by Frederick Law Olmsted


\item The Plateau: colourful row houses and lively cafés


\end{enumerate}

Local specialties worth trying:

\begin{itemize}
\item Poutine


\item Smoked meat sandwiches


\item Bagels from \href{https://www.stviateurbagel.com}{St-Viateur}


\end{itemize}

The city's motto is \texttt{Concordia Salus}\footnotemark[1], meaning ``well-being through harmony.''

As the Montréal tourism board puts it:\\
``\emph{There is no place quite like it.}''



\protect\includegraphics{/assets/placeholder.png}\subsection*{Code chunks}
\label{code-chunks}

Code chunks are fenced with triple backticks and a \texttt{\{r\}} header. Options go inside the chunk as \texttt{\#\textbar{}} pipe comments. See the \href{/code/chunks.html}{Chunks} vignette for the full list of chunk options.

\footnotetext[1]{The motto was adopted in 1833 and appears on the city's coat of arms.

}


\begin{srccode}
\begin{Verbatim}[commandchars=\\\{\}]
\textcolor{hl1}{```\{r, chunk}\textcolor{hl2}{-}\textcolor{hl3}{1}\textcolor{hl1}{\}}
\textcolor{hl4}{\#| echo = fenced
}
\textcolor{hl1}{x }\textcolor{hl2}{<-}\textcolor{hl1}{ }\textcolor{hl3}{1}\textcolor{hl2}{:}\textcolor{hl3}{10}
\textcolor{hl3}{mean}\textcolor{hl1}{(x)}
\textcolor{hl1}{```}
\end{Verbatim}
\end{srccode}

\begin{outcode}
\begin{verbatim}
> [1] 5.5
\end{verbatim}
\end{outcode}


\begin{figure}
\centering
\includegraphics[width=0.70\textwidth]{_calepin/files/chunk-1-1.pdf}\label{chunk-1}
\end{figure}

\subsection*{Inline code}
\label{inline-code}

R expressions can be evaluated inline using back ticks and curly braces. See the \href{/code/inline.html}{Inline} vignette for more details.




\begin{srccode}
\begin{Verbatim}[commandchars=\\\{\}]
\textcolor{hl1}{`\{r\} round(pi, 4)`}
\end{Verbatim}
\end{srccode}

The numbers in this sentence are computed dynamically: The mean of 1:10 is 5.5 and pi is approximately 3.14.

\subsection*{Figures}
\label{figures}




\begin{srccode}
\begin{Verbatim}[commandchars=\\\{\}]
\textcolor{hl3}{plot}\textcolor{hl1}{(}
\textcolor{hl1}{    }\textcolor{hl3}{1}\textcolor{hl2}{:}\textcolor{hl3}{10}\textcolor{hl1}{,}
\textcolor{hl1}{    (}\textcolor{hl3}{1}\textcolor{hl2}{:}\textcolor{hl3}{10}\textcolor{hl1}{)}\textcolor{hl2}{\^{}}\textcolor{hl3}{2}\textcolor{hl1}{,}
\textcolor{hl1}{    }\textcolor{hl3}{main}\textcolor{hl1}{ }\textcolor{hl2}{=}\textcolor{hl1}{ }\textcolor{hl5}{"Square Numbers"}\textcolor{hl1}{,}
\textcolor{hl1}{    }\textcolor{hl3}{xlab}\textcolor{hl1}{ }\textcolor{hl2}{=}\textcolor{hl1}{ }\textcolor{hl5}{"x"}\textcolor{hl1}{,}
\textcolor{hl1}{    }\textcolor{hl3}{ylab}\textcolor{hl1}{ }\textcolor{hl2}{=}\textcolor{hl1}{ }\textcolor{hl5}{"x\^{}2"}
\textcolor{hl1}{)}
\end{Verbatim}
\end{srccode}


\begin{figure}
\centering
\includegraphics[width=0.70\textwidth]{_calepin/files/fig-scatter-1.pdf}
\caption{A simple scatter plot}\label{fig-scatter}
\end{figure}

\subsection*{Cross-references}
\label{cross-references}

\emph{Calepin} resolves references to figures, sections, and theorem environments.

\begin{itemize}
\item \hyperref[fig-scatter]{Figure 1} renders as a linked reference


\item Section \ref{sec-figures} links to the Figures section


\item \hyperref[fig-scatter]{1} suppresses the type prefix


\end{itemize}

Sections get automatic ids from their heading text.

\subsection*{Math}
\label{math}

Inline math: $a^2 + b^2 = c^2$. Display math:

$$
\int_0^1 x^2 \, dx = \frac{1}{3}
$$



\par\medskip\noindent\textbf{Theorem 1.}\label{thm-pythagoras} \itshape{} In a right triangle, the square of the hypotenuse equals the sum of the squares of the other two sides: $a^2 + b^2 = c^2$.

\upshape\medskip

\par\medskip\noindent\textit{.} Let a right triangle have legs $a$, $b$ and hypotenuse $c$. By constructing squares on each side and comparing areas, we obtain $a^2 + b^2 = c^2$.

 \hfill$\square$\medskip

By \hyperref[thm-pythagoras]{Theorem 1}, the relationship is fundamental to Euclidean geometry.

\subsection*{Callouts}
\label{callouts}

Callout divs highlight important information. Five types are available: \texttt{callout-note}, \texttt{callout-warning}, \texttt{callout-tip}, \texttt{callout-caution}, and \texttt{callout-important}. They can be nested.




\begin{tcolorbox}[enhanced, breakable, colback=blue!5, colframe=blue!40,
  left=4pt, right=4pt, top=4pt, bottom=4pt,
  boxrule=0pt, leftrule=3pt, arc=0pt, outer arc=0pt,
  title=Note, fonttitle=\bfseries]
This is a note callout. Use it for supplementary information.




\begin{tcolorbox}[enhanced, breakable, colback=yellow!10, colframe=yellow!60!black,
  left=4pt, right=4pt, top=4pt, bottom=4pt,
  boxrule=0pt, leftrule=3pt, arc=0pt, outer arc=0pt,
  title=Warning, fonttitle=\bfseries]
Callouts can be nested. This warning appears inside the note above.



\end{tcolorbox}

\end{tcolorbox}

\subsection*{Conditional content}
\label{conditional-content}

Fenced divs with \texttt{.content-visible} or \texttt{.content-hidden} classes control which content appears in each output format.



Format aliases are supported: \texttt{latex} and \texttt{pdf} both match LaTeX output, \texttt{markdown} and \texttt{md} both match Markdown.

\subsection*{Raw blocks}
\label{raw-blocks}

Raw blocks inject format-specific markup that passes through verbatim when the output format matches and is dropped otherwise.



\subsection*{Tabsets}
\label{tabsets}

Tabsets organize content into switchable tabs. In HTML, readers click to switch tabs. In LaTeX and Typst, tabs render as regular sections.



\subsubsection*{R}
\label{r}




\begin{srccode}
\begin{Verbatim}[commandchars=\\\{\}]
\textcolor{hl1}{x }\textcolor{hl2}{<-}\textcolor{hl1}{ }\textcolor{hl3}{c}\textcolor{hl1}{(}\textcolor{hl3}{1}\textcolor{hl1}{, }\textcolor{hl3}{2}\textcolor{hl1}{, }\textcolor{hl3}{3}\textcolor{hl1}{)}
\textcolor{hl3}{mean}\textcolor{hl1}{(x)}
\end{Verbatim}
\end{srccode}

\begin{outcode}
\begin{verbatim}
> [1] 2
\end{verbatim}
\end{outcode}


\begin{figure}
\centering
\includegraphics[width=0.70\textwidth]{_calepin/files/chunk-3-1.pdf}\label{chunk-3}
\end{figure}

\subsubsection*{Python}
\label{python}




\begin{srccode}
\begin{Verbatim}[commandchars=\\\{\}]
\textcolor{hl1}{x }\textcolor{hl2}{=}\textcolor{hl1}{ [}\textcolor{hl3}{1}\textcolor{hl1}{, }\textcolor{hl3}{2}\textcolor{hl1}{, }\textcolor{hl3}{3}\textcolor{hl1}{]}
\textcolor{hl3}{sum}\textcolor{hl1}{(x) }\textcolor{hl2}{/}\textcolor{hl1}{ }\textcolor{hl3}{len}\textcolor{hl1}{(x)}
\end{Verbatim}
\end{srccode}

\subsubsection*{Julia}
\label{julia}




\begin{srccode}
\begin{Verbatim}[commandchars=\\\{\}]
\textcolor{hl1}{x = [1, 2, 3]
}
\textcolor{hl1}{sum(x) / length(x)}
\end{Verbatim}
\end{srccode}

Tabsets with the same \texttt{group} attribute switch together:



\subsubsection*{R}
\label{r}

R version of example 2.

\subsubsection*{Python}
\label{python}

Python version of example 2.



\subsubsection*{R}
\label{r}

R content stays synced with the tabset above.

\subsubsection*{Python}
\label{python}

Python content stays synced too.



\subsection*{Line blocks}
\label{line-blocks}

Line blocks preserve line breaks and leading spaces. Prefix each line with \texttt{\textbar{}}.

The limerick packs laughs anatomical\\
Into space that is quite economical.\\
~~~But the good ones I've seen\\
~~~So seldom are clean\\
And the clean ones so seldom are comical.

\subsection*{Footnotes}
\label{footnotes}

There are two ways to create footnotes. Named references use \texttt{{[}\^{}label{]}} at the point of reference and \texttt{{[}\^{}label{]}: text} for the definition elsewhere in the document. Inline footnotes use \texttt{\^{}{[}text{]}} directly without a separate definition.

Named footnote\footnotemark[2] and inline footnote\footnotemark[3].

\subsection*{Numbered sections}
\label{numbered-sections}

Add \texttt{number-sections: true} to the YAML front matter to automatically number all section headings.

\subsection*{The \texttt{.hidden} div}
\label{the-div}

The \texttt{.hidden} div executes its content but produces no output.

\footnotetext[2]{This is defined separately.

}\footnotetext[3]{This is defined right here.

}

The value computed silently: \_\_CALEPIN\_1aaf\_2\_\_\_ERROR:object `secret' not found.






\end{document}